
% \chapter{Miscellany}
\chapter{The \texorpdfstring{\mpiarg{Info}}{Info} Object}
\mpitermtitleindex{info object}
\label{sec:misc-2}
\label{chap:misc-2}

%%%%%%%%%%%%%%%%%%%%%%%%%%%%%%%%%%%%%%%%%%%%%%%%%%%%%%%%%%%%%%%%%%%%%%

%This chapter contains topics that do not fit conveniently into other chapters.


% Moved here from Dynamic
\label{subsec:info}
\label{chap:info}

Many of the procedures in 
\MPI/ 
take an argument
\mpiarg{info}. \mpiarg{info} is an opaque object with a handle of type
\cdeclmainindex{MPI\_Info}\cdeclmainindex{TYPE(MPI\_Info)}%
\mpictype{MPI\_Info} in C and Fortran with the \code{mpi\_f08} module, 
and 
\ftype{INTEGER} in Fortran with the \code{mpi} module or the (deprecated) \code{mpif.h} include file. 
It stores an unordered set of
(\mpiarg{key},\mpiarg{value}) pairs (both \mpiarg{key} and
\mpiarg{value} are strings). 
A key 
can
have only one
value. \MPI/ reserves several keys and requires that if an
implementation uses a reserved key, it must provide the specified
functionality.  An implementation is not required to support these
keys and may support any others not reserved by \MPI/.

Some info hints allow the \MPI/ library to
restrict its support for certain operations
in order to improve performance or resource utilization.  If an
application provides such an info hint, it must be compatible with
any changes in the behavior of the \MPI/ library that are allowed by the info
hint.


An implementation must support info objects as caches for arbitrary (\mpiarg{key},\mpiarg{value}) pairs, regardless of whether it recognizes the key. Each procedure that
takes hints in the form of an \mpictype{MPI\_Info} must be prepared to ignore any key it
does not recognize. This description of info objects does not attempt to
define how a particular procedure should react if it recognizes a key but not the
associated value. 
\mpifunc{MPI\_INFO\_GET\_NKEYS}, \mpifunc{MPI\_INFO\_GET\_NTHKEY},
and \mpifunc{MPI\_INFO\_GET\_STRING}
must retain all (\mpiarg{key},\mpiarg{value})
pairs so that layered functionality can also use the \mpiarg{Info} object. 

\begin{sloppy}
Keys have an implementation-defined maximum length of
\mpiconstmain{MPI\_MAX\_INFO\_KEY}-1,
where \mpiconst{MPI\_MAX\_INFO\_KEY} is at least 33 and at most 256.
Values have an implementation-defined maximum length of
\mpiconstmain{MPI\_MAX\_INFO\_VAL}.  
In Fortran, leading and trailing spaces are stripped from
both.  Returned values will never be larger than these maximum
lengths.
Both \mpiarg{key} and \mpiarg{value} are case sensitive.

\end{sloppy}

\begin{rationale}
Keys have a maximum length because the set of known keys will always
be finite and known to the implementation and because there is no
reason for keys to be complex.  The small maximum size allows
applications to declare keys of size \mpiconst{MPI\_MAX\_INFO\_KEY}.
The limitation on value sizes is so that an implementation is not
forced to deal with arbitrarily long
strings.
\end{rationale}
\begin{users}
\mpiconst{MPI\_MAX\_INFO\_VAL} might be very large, so it
might not be wise to declare a string of that size.
\end{users}

When \mpiarg{info} is used as an argument to any \MPI/ procedure, it
is interpreted before that procedure returns, so that it may
be read, modified, or freed immediately after return.
Changes to an info object after return from a procedure do not affect
that interpretation.

\begin{rationale}
Prior to \mpiivdoto/, the above statement was restricted to nonblocking \MPI/ procedures.
For simplicity this restriction was removed, as it currently applies to all \MPI/
procedures that use \mpiarg{info} arguments. Note, this has to be revisited for new
procedures added in the future, e.g., for future procedures that could return an \mpiarg{info}
argument to be filled in after the return from the procedure.
\end{rationale}

When the descriptions refer to a key or value as being a boolean, an
integer, or a list, they mean the string representation of these
types.  An implementation may define its own rules for how info value
strings are converted to other types, but to ensure portability, every
implementation must support the following representations.  Valid
values for a boolean must include the strings \infovalmain{true} and \infovalmain{false}
(all lowercase).  For integers, valid values must include string
representations of decimal values of integers that are within the
range of a standard integer type in the program.  (However it is
possible that not every integer is a valid value for a given
key.)  On positive numbers, $+$ signs are optional.  No space may
appear between a $+$ or $-$ sign and the leading digit of a number.  For
comma separated lists, the string must contain valid elements
separated by commas.  Leading and trailing spaces are stripped
automatically from the types of info values described above and for
each element of a comma separated list.  These rules apply to all info
values of these types.  Implementations are free to specify a
different interpretation for values of other info keys.

\cdeclindex{MPI\_Info}%
\begin{funcdef}{MPI\_INFO\_CREATE}{(\mbox{info})}
\funcarg{\OUT}{info}{info object created (handle)}
\end{funcdef}
\mpibindmain{MPI\_Info\_create}{(\gb{}MPI\_Info~*info)}
\mpifnewbindmain{MPI\_Info\_create}{(info, ierror) \fargs TYPE(MPI\_Info), INTENT(OUT)\ ::\ \mbox{info}\fnfinalarg{}INTEGER, OPTIONAL, INTENT(OUT)\ ::\ \mbox{ierror}}
\mpifbindmain{MPI\_INFO\_CREATE}{(INFO, IERROR) \fargs \fnfinalarg{}INTEGER \mbox{INFO}, \mbox{IERROR}}

\mpifunc{MPI\_INFO\_CREATE} creates a new info object. The newly created object contains
no key/value pairs. 

\cdeclindex{MPI\_Info}%
\begin{funcdef}{MPI\_INFO\_SET}{(\mbox{info}, \mbox{key}, \mbox{value})}
\funcarg{\INOUT}{info}{info object (handle)}
\funcarg{\IN}{key}{key (string)}
\funcarg{\IN}{value}{value (string)}
\end{funcdef}
\mpibindmain{MPI\_Info\_set}{(\gb{}MPI\_Info~info, const~char~*key, const~char~*value)}
\mpifnewbindmain{MPI\_Info\_set}{(info, key, value, ierror) \fargs TYPE(MPI\_Info), INTENT(IN)\ ::\ \mbox{info}\fnarg{}CHARACTER(LEN=*), INTENT(IN)\ ::\ \mbox{key}, \mbox{value}\fnfinalarg{}INTEGER, OPTIONAL, INTENT(OUT)\ ::\ \mbox{ierror}}
\mpifbindmain{MPI\_INFO\_SET}{(INFO, KEY, VALUE, IERROR) \fargs INTEGER \mbox{INFO}, \mbox{IERROR}\fnfinalarg{}CHARACTER*(*) \mbox{KEY}, \mbox{VALUE}}

\mpifunc{MPI\_INFO\_SET} adds the (\mpiarg{key},\mpiarg{value}) pair to \mpiarg{info}, and  overrides the value if a value for the same key was previously set. 
\mpiarg{key} and \mpiarg{value} are null-terminated strings in C. 
In Fortran, leading and trailing spaces in \mpiarg{key} and \mpiarg{value} are stripped. 
If either \mpiarg{key} or \mpiarg{value} are longer than 
the respective maximum length, the call raises an error of class \errormain{MPI\_ERR\_INFO\_KEY} or
\errormain{MPI\_ERR\_INFO\_VALUE}, respectively.

\cdeclindex{MPI\_Info}%
\begin{funcdef}{MPI\_INFO\_DELETE}{(\mbox{info}, \mbox{key})}
\funcarg{\INOUT}{info}{info object (handle)}
\funcarg{\IN}{key}{key (string)}
\end{funcdef}
\mpibindmain{MPI\_Info\_delete}{(\gb{}MPI\_Info~info, const~char~*key)}
\mpifnewbindmain{MPI\_Info\_delete}{(info, key, ierror) \fargs TYPE(MPI\_Info), INTENT(IN)\ ::\ \mbox{info}\fnarg{}CHARACTER(LEN=*), INTENT(IN)\ ::\ \mbox{key}\fnfinalarg{}INTEGER, OPTIONAL, INTENT(OUT)\ ::\ \mbox{ierror}}
\mpifbindmain{MPI\_INFO\_DELETE}{(INFO, KEY, IERROR) \fargs INTEGER \mbox{INFO}, \mbox{IERROR}\fnfinalarg{}CHARACTER*(*) \mbox{KEY}}

\mpifunc{MPI\_INFO\_DELETE} deletes a (\mpiarg{key},\mpiarg{value}) pair from \mpiarg{info}.
If \mpiarg{key} is not defined in \mpiarg{info}, the call
raises an error of class \errormain{MPI\_ERR\_INFO\_NOKEY}.

\cdeclindex{MPI\_Info}%
\begin{funcdef}{MPI\_INFO\_GET\_STRING}{(\mbox{info}, \mbox{key}, \mbox{buflen}, \mbox{value}, \mbox{flag})}
\funcarg{\IN}{info}{info object (handle)}
\funcarg{\IN}{key}{key (string)}
\funcarg{\INOUT}{buflen}{length of buffer (integer)}
\funcarg{\OUT}{value}{value (string)}
\funcarg{\OUT}{flag}{\mpicode{true} if \mpiarg{key} is defined, \mpicode{false} otherwise (logical)}
\end{funcdef}
\mpibindmain{MPI\_Info\_get\_string}{(\gb{}MPI\_Info~info, const~char~*key, int~*buflen, char~*value, int~*flag)}
\mpifnewbindmain{MPI\_Info\_get\_string}{(info, key, buflen, value, flag, ierror) \fargs TYPE(MPI\_Info), INTENT(IN)\ ::\ \mbox{info}\fnarg{}CHARACTER(LEN=*), INTENT(IN)\ ::\ \mbox{key}\fnarg{}INTEGER, INTENT(INOUT)\ ::\ \mbox{buflen}\fnarg{}CHARACTER(LEN=*), INTENT(OUT)\ ::\ \mbox{value}\fnarg{}LOGICAL, INTENT(OUT)\ ::\ \mbox{flag}\fnfinalarg{}INTEGER, OPTIONAL, INTENT(OUT)\ ::\ \mbox{ierror}}
\mpifbindmain{MPI\_INFO\_GET\_STRING}{(INFO, KEY, BUFLEN, VALUE, FLAG, IERROR) \fargs INTEGER \mbox{INFO}, \mbox{BUFLEN}, \mbox{IERROR}\fnarg{}CHARACTER*(*) \mbox{KEY}, \mbox{VALUE}\fnfinalarg{}LOGICAL \mbox{FLAG}}



This procedure retrieves the value associated with \mpiarg{key} from \mpiarg{info}, if any.
If such a key exists in \mpiarg{info}, it sets \mpiarg{flag} to \mpicode{true}
and returns the value in \mpiarg{value},
otherwise it sets \mpiarg{flag} to \mpicode{false} and leaves
\mpiarg{value} unchanged. 
\mpiarg{buflen} on input is the size of the provided buffer, \mpiarg{value}, for the output of \mpiarg{buflen} it is the size of the buffer needed to store the
value string. If the \mpiarg{buflen} passed into the procedure is less than the
actual size needed to store the value string (including null terminator in C), the value is
truncated. On return, the value of \mpiarg{buflen} will be set to the required
buffer size to hold the value string. If \mpiarg{buflen} is set to \mpicode{0}, \mpiarg{value} is not changed. 
In C, \mpiarg{buflen} 
includes the required space for the null terminator. In C,
this procedure returns a null terminated string in all cases where
the \mpiarg{buflen} input value is greater than \mpicode{0}.

If \mpiarg{key} is larger than \mpiconst{MPI\_MAX\_INFO\_KEY}, 
the call is erroneous. 

\begin{users}
The \mpifunc{MPI\_INFO\_GET\_STRING} procedure can be used to obtain the size of
the required buffer for a value string by setting the \mpiarg{buflen} to \mpicode{0}. The returned
\mpiarg{buflen} can then be used to allocate memory before calling
\mpifunc{MPI\_INFO\_GET\_STRING} again to obtain the value string.
\end{users}

\cdeclindex{MPI\_Info}%
\begin{funcdef}{MPI\_INFO\_GET\_NKEYS}{(\mbox{info}, \mbox{nkeys})}
\funcarg{\IN}{info}{info object (handle)}
\funcarg{\OUT}{nkeys}{number of defined keys (integer)}
\end{funcdef}
\mpibindmain{MPI\_Info\_get\_nkeys}{(\gb{}MPI\_Info~info, int~*nkeys)}
\mpifnewbindmain{MPI\_Info\_get\_nkeys}{(info, nkeys, ierror) \fargs TYPE(MPI\_Info), INTENT(IN)\ ::\ \mbox{info}\fnarg{}INTEGER, INTENT(OUT)\ ::\ \mbox{nkeys}\fnfinalarg{}INTEGER, OPTIONAL, INTENT(OUT)\ ::\ \mbox{ierror}}
\mpifbindmain{MPI\_INFO\_GET\_NKEYS}{(INFO, NKEYS, IERROR) \fargs \fnfinalarg{}INTEGER \mbox{INFO}, \mbox{NKEYS}, \mbox{IERROR}}

\mpifunc{MPI\_INFO\_GET\_NKEYS} returns the number of currently defined keys in \mpiarg{info}.


\cdeclindex{MPI\_Info}%
\begin{funcdef}{MPI\_INFO\_GET\_NTHKEY}{(\mbox{info}, \mbox{n}, \mbox{key})}
\funcarg{\IN}{info}{info object (handle)}
\funcarg{\IN}{n}{key number (integer)}
\funcarg{\OUT}{key}{key (string)}
\end{funcdef}
\mpibindmain{MPI\_Info\_get\_nthkey}{(\gb{}MPI\_Info~info, int~n, char~*key)}
\mpifnewbindmain{MPI\_Info\_get\_nthkey}{(info, n, key, ierror) \fargs TYPE(MPI\_Info), INTENT(IN)\ ::\ \mbox{info}\fnarg{}INTEGER, INTENT(IN)\ ::\ \mbox{n}\fnarg{}CHARACTER(LEN=*), INTENT(OUT)\ ::\ \mbox{key}\fnfinalarg{}INTEGER, OPTIONAL, INTENT(OUT)\ ::\ \mbox{ierror}}
\mpifbindmain{MPI\_INFO\_GET\_NTHKEY}{(INFO, N, KEY, IERROR) \fargs INTEGER \mbox{INFO}, \mbox{N}, \mbox{IERROR}\fnfinalarg{}CHARACTER*(*) \mbox{KEY}}

This procedure returns the \mpiarg{n}th defined key in \mpiarg{info}.
Keys are numbered $0 \ldots N-1$ where $N$ is the
value returned by \mpifunc{MPI\_INFO\_GET\_NKEYS}. 
All keys between $0$ and $N-1$ are guaranteed to 
be defined. The number of a given key does not change
as long as \mpiarg{info} is not modified with 
\mpifunc{MPI\_INFO\_SET} or \mpifunc{MPI\_INFO\_DELETE}.


\cdeclindex{MPI\_Info}%
\begin{funcdef}{MPI\_INFO\_DUP}{(\mbox{info}, \mbox{newinfo})}
\funcarg{\IN}{info}{info object (handle)}
\funcarg{\OUT}{newinfo}{info object created (handle)}
\end{funcdef}
\mpibindmain{MPI\_Info\_dup}{(\gb{}MPI\_Info~info, MPI\_Info~*newinfo)}
\mpifnewbindmain{MPI\_Info\_dup}{(info, newinfo, ierror) \fargs TYPE(MPI\_Info), INTENT(IN)\ ::\ \mbox{info}\fnarg{}TYPE(MPI\_Info), INTENT(OUT)\ ::\ \mbox{newinfo}\fnfinalarg{}INTEGER, OPTIONAL, INTENT(OUT)\ ::\ \mbox{ierror}}
\mpifbindmain{MPI\_INFO\_DUP}{(INFO, NEWINFO, IERROR) \fargs \fnfinalarg{}INTEGER \mbox{INFO}, \mbox{NEWINFO}, \mbox{IERROR}}

\mpifunc{MPI\_INFO\_DUP} duplicates an existing info object, creating a new
object, with the same (\mpiarg{key},\mpiarg{value}) pairs and the same ordering of keys.

\cdeclindex{MPI\_Info}%
\begin{funcdef}{MPI\_INFO\_FREE}{(\mbox{info})}
\funcarg{\INOUT}{info}{info object (handle)}
\end{funcdef}
\mpibindmain{MPI\_Info\_free}{(\gb{}MPI\_Info~*info)}
\mpifnewbindmain{MPI\_Info\_free}{(info, ierror) \fargs TYPE(MPI\_Info), INTENT(INOUT)\ ::\ \mbox{info}\fnfinalarg{}INTEGER, OPTIONAL, INTENT(OUT)\ ::\ \mbox{ierror}}
\mpifbindmain{MPI\_INFO\_FREE}{(INFO, IERROR) \fargs \fnfinalarg{}INTEGER \mbox{INFO}, \mbox{IERROR}}

This procedure frees \mpiarg{info} and sets it to \mpiconstmain{MPI\_INFO\_NULL}.

\begin{funcdef}{MPI\_INFO\_CREATE\_ENV}{(\mbox{info})}
\funcarg{\OUT}{info}{info object (handle)}
\end{funcdef}
\mpibindmain{MPI\_Info\_create\_env}{(\gb{}int~argc, char~*argv[], MPI\_Info~*info)}
\mpifnewbindmain{MPI\_Info\_create\_env}{(info, ierror) \fargs TYPE(MPI\_Info), INTENT(OUT)\ ::\ \mbox{info}\fnfinalarg{}INTEGER, OPTIONAL, INTENT(OUT)\ ::\ \mbox{ierror}}
\mpifbindmain{MPI\_INFO\_CREATE\_ENV}{(INFO, IERROR) \fargs \fnfinalarg{}INTEGER \mbox{INFO}, \mbox{IERROR}}

This procedure creates an output object \mpiarg{info} with the same construction as
\mpiconst{MPI\_INFO\_ENV} as created during \mpifunc{MPI\_INIT} or \mpifunc{MPI\_INIT\_THREAD} when the
same arguments are used.
This construction is described in Section~\ref{sec:inquiry-startup};
however, this procedure can be called when not using the \wpm{}, e.g., when using the \spm{}.
This object is not a direct
copy or alias of the \mpiconst{MPI\_INFO\_ENV} object and could contain different values based on the
input arguments and other sources.  Multiple calls to this procedure that are given the same input
arguments will produce \mpiarg{info} objects consistent with the definition of
\mpiconst{MPI\_INFO\_ENV}.  The version for ISO C accepts the \mpiarg{argc} and \mpiarg{argv} that are
provided by the arguments to \code{main} or \code{0} for \mpiarg{argc} and \code{NULL} for
\mpiarg{argv}.  The user is responsible for freeing the \mpiarg{info} object via
\mpifunc{MPI\_INFO\_FREE}.  This procedure is local.

This procedure must always be thread-safe, as defined in Section~\ref{sec:dynamic:threads}.  It is
one of the few procedures that may be called before \MPI/ is initialized or after \MPI/ is finalized.

\begin{users}

    In some circumstances (e.g., when passing \code{0} to \mpiarg{argc} and \code{NULL} to
    \mpiarg{argv} in C or in Fortran where such arguments do not exist), the \mpiarg{info} object
    may not be populated or may be populated incompletely because this procedure is local and the
    implementation may not be able to determine the correct values. Note that this could result in
    different values in the resulting \mpiarg{info} object at different \MPI/ processes.

\end{users}

